\section{A contrasting critical law for overlapping cycles}\label{sec:overlap}

Keep $V=\mathbb F_q^{2n}$, $N=q^{2n}$ and the nonzero-vector sampling model.
Let $2\le a\le2n$. Assign independent uniform nonzero vectors
$X_1,\ldots,X_a$ to the left part of $K_{a,m}$ and
$Z_1,\ldots,Z_m$ to the right part. Orient all edges from left to right and
observe the columns
\[
 Y_j=(\langle X_1,Z_j\rangle,\ldots,
       \langle X_a,Z_j\rangle)\in\mathbb F_q^a.
\]
Let $P_{a,m,N}$ be the actual law and let
$Q_{a,m,N}=(p_N^{\otimes a})^{\otimes m}$ be the law with independent
edges and the same one-edge marginal.

Define a probability measure on $\mathbb F_q^a$ by
\[
 \rho_{a,N}(y)=
 \begin{cases}
 (N-q^a)/(q^a(N-1)),&y=0,\\
 N/(q^a(N-1)),&y\ne0.
 \end{cases}
\]
Its zero mass is allowed to vanish at the finite endpoint $a=2n$.
Put
\[
 \alpha_{a,N}=\prod_{i=0}^{a-1}\frac{N-q^i}{N-1},\qquad
 \delta_{a,N}=1-\alpha_{a,N}.
\]

\begin{proposition}[Exact independent-anchor reduction]\label{prop:anchors}
The event $\mathcal E$ that the $a$ anchor vectors are linearly independent
has probability $\alpha_{a,N}$. Conditional on $\mathcal E$, the columns
$Y_1,\ldots,Y_m$ are independent with common law $\rho_{a,N}$.
Consequently,
\[
 d_{\mathrm{TV}}(P_{a,m,N},\rho_{a,N}^{\otimes m})
 \le\delta_{a,N}
 \le\frac{(q^a-1)/(q-1)-a}{N-1}.
\tag{A1}
\]
Both inequalities hold for every $m$.

\end{proposition}
\begin{proof}
After $i$ independent nonzero anchors, the span contains $q^i-1$ of the
$N-1$ possible next vectors. Multiplying the successive probabilities
proves the formula for $\alpha_{a,N}$. The elementary product inequality
$1-\prod_i(1-b_i)\le\sum_i b_i$ with
$b_i=(q^i-1)/(N-1)$ gives the second bound.

For each fixed independent anchor tuple, nondegeneracy of the alternating
form makes the linear map
\[
 T_X:V\longrightarrow\mathbb F_q^a,
 \qquad z\longmapsto(\langle X_i,z\rangle)_{i=1}^a
\]
surjective. Every fibre has $N/q^a$ elements. Removing the zero vector
removes one point from the zero fibre only, giving $\rho_{a,N}$.
The $Z_j$ are independent, and this conditional column law does not depend
on the particular independent anchor tuple. Integrating over all such
tuples therefore preserves the product law. The mixture decomposition
\[
 P_{a,m,N}=\alpha_{a,N}\rho_{a,N}^{\otimes m}
             +\delta_{a,N}\mathcal D_{a,m,N}
\]
for a probability measure $\mathcal D_{a,m,N}$ proves the first bound.
\end{proof}

\begin{lemma}[Conditional column information]\label{lem:anchor-info}
For fixed $a,q$ and $n\to\infty$, uniformly on $y\in\mathbb F_q^a$,
\[
 \frac{\rho_{a,N}(y)}{p_N^{\otimes a}(y)}
 =1+\frac{h_a(y)}N+O_{a,q}(N^{-2}),
 \qquad
 h_a(y)=\sum_{i=1}^a\kappa(y_i)-\bigl(q^a\mathbf1_{y=0}-1\bigr).
\tag{A2}
\]
Under independent uniform coordinates, this score has mean zero and
\[
 \mathbb E h_a^2=J_{a,q}:=q^a-1-a(q-1)>0.
\tag{A3}
\]
In particular,
\[
 \chi^2(\rho_{a,N}\Vert p_N^{\otimes a})
 =J_{a,q}N^{-2}+O_{a,q}(N^{-3}).
\tag{A4}
\]

\end{lemma}
\begin{proof}
For $z=\#\{i:y_i=0\}$ and $y\ne0$, the likelihood is exactly
\[
 (1-N^{-1})^{a-1}(1-q/N)^{-z}.
\]
For $y=0$ it is
\[
 (1-q^a/N)(1-N^{-1})^{a-1}(1-q/N)^{-a}.
\]
Expanding these finite rational functions proves (A2).
Under uniform coordinates, the variables $\kappa(y_i)$ are independent,
have mean zero, and have second moment $q-1$. Moreover,
\[
 q^a\mathbf1_{y=0}=\prod_{i=1}^a(1+\kappa(y_i)),
 \qquad
 h_a(y)=-\sum_{\substack{S\subseteq\{1,\ldots,a\}\\|S|\ge2}}
               \prod_{i\in S}\kappa(y_i).
\]
The summands are pairwise orthogonal in $L^2$ of uniform product measure.
Their squared norms sum to
$\sum_{s=2}^a\binom as(q-1)^s=q^a-1-a(q-1)$, proving (A3).
The actual reference measure $p_N^{\otimes a}$ differs pointwise from the
uniform measure by $O_{a,q}(N^{-1})$, and (A2) then proves (A4).

For finite calculations there is also the exact identity
\[
 1+\chi^2(\rho_{a,N}\Vert p_N^{\otimes a})
 =\frac{(N-1)^{a-2}}{(N-q)^a}
   \left[N^{2-a}(N-q+1)^a+q^a-2N\right].
\tag{A5}
\]
To verify it, write
$\rho_{a,N}(y)=c(1-q^aN^{-1}\mathbf1_{y=0})$ with
$c=N/[q^a(N-1)]$, and sum its squared density over the product reference.
The two necessary sums are
\[
 \sum_{b\in\mathbb F_q}\frac1{p_N(b)}
 =\frac{q^2(N-1)(N-q+1)}{N(N-q)},\qquad
 \frac1{p_N(0)^a}=\left(\frac{q(N-1)}{N-q}\right)^a.
\]
Substitution and cancellation give (A5).
\end{proof}

\begin{theorem}[The fixed-anchor critical profile]\label{thm:anchors}
Fix $q$ and $a\ge2$. Let $n\to\infty$ and let $m=m_n$ be arbitrary
positive integers. If
\[
 \frac{mJ_{a,q}}{N^2}\longrightarrow\tau\in[0,\infty],
\]
then
\[
 d_{\mathrm{TV}}(P_{a,m,N},Q_{a,m,N})
 \longrightarrow 2\Phi(\sqrt\tau/2)-1,
\tag{A6}
\]
with value $1$ at $\tau=\infty$.
\end{theorem}
\begin{proof}
Indeed, apply \cref{lem:lan} to
$\rho_{a,N}^{\otimes m}$ and $(p_N^{\otimes a})^{\otimes m}$, using
(A2)--(A4). The difference between their distance and the actual distance
is at most $\delta_{a,N}=O_{a,q}(N^{-1})$ by (A1) and the triangle
inequality. This error is uniform in $m$.
\end{proof}

\subsection{Many overlapping four-cycles with vanishing total variation}
Set $a=2$. Now
\[
 \delta_{2,N}=\frac{q-1}{N-1},\qquad
 \chi^2(\rho_{2,N}\Vert p_N^{\otimes2})
      =\frac{(q-1)^2}{(N-q)^2}.
\]
The exact product chi-square identity and Cauchy--Schwarz give
\[
 d_{\mathrm{TV}}(P_{2,m,N},Q_{2,m,N})
 \le\frac{q-1}{N-1}
 +\frac12\sqrt{\left(1+\frac{(q-1)^2}{(N-q)^2}\right)^m-1}.
\tag{A7}
\]
The graph $K_{2,m}$ has exactly $\binom m2$ simple four-cycles.
Each four-cycle marginal has chi-square divergence
\[
 v_{4,N}=(q-1)^2\left((N-q)^{-2}+(q-2)(N-1)^{-2}\right)
       \sim(q-1)^3N^{-2}.
\]
Choose any fixed $1<\beta<2$ and $m=\lfloor N^\beta\rfloor$.
Then (A7) tends to zero, while
\[
 \binom m2 v_{4,N}\sim\frac{(q-1)^3}{2}N^{2\beta-2}
 \longrightarrow\infty.
\tag{A8}
\]
Thus summing the information of overlapping cycle marginals does not give
the information of the whole observation. The exact cactus factorization
is an essential hypothesis of the cycle-count criterion.

\subsection{The full chi-square divergence can diverge exponentially too}
Fix $c\in\mathbb F_q^\times$, and let $E_c$ be the data event
$Y_{2j}=cY_{1j}$ for every $1\le j\le m$.
The latent event $X_2=cX_1$ has probability $1/(N-1)$ and forces $E_c$,
so $P_{2,m,N}(E_c)\ge1/(N-1)$.
Under the independent-edge law,
\[
 Q_{2,m,N}(E_c)=s_N^m,\qquad
 s_N=p_N(0)^2+(q-1)p_N(1)^2
    =\frac1q+\frac{q-1}{q(N-1)^2}.
\]
Since $1+\chi^2(P\Vert Q)=\sum_y P(y)^2/Q(y)$, applying
Cauchy--Schwarz on $E_c$ gives the exact lower bound
\[
 1+\chi^2(P_{2,m,N}\Vert Q_{2,m,N})
 \ge\frac1{(N-1)^2s_N^m}.
\tag{A9}
\]
For the same $m=\lfloor N^\beta\rfloor$, $1<\beta<2$,
\[
 \log(1+\chi^2(P_{2,m,N}\Vert Q_{2,m,N}))
 \ge m\log q-2\log(N-1)-O_q(m/N^2)
 \longrightarrow\infty,
\]
even though the total variation distance tends to zero.
The rare dependent-anchor configurations therefore have negligible total
probability and arbitrarily large likelihood ratios. This explicitly
explains why a direct second-moment bound for the original mixture may be
very poor, whereas conditioning on anchor independence identifies its
correct testing profile.
